\documentclass[a4paper,12pt]{article}
\usepackage[spanish]{babel}
\usepackage[utf8]{inputenc}
\usepackage[T1]{fontenc}
\usepackage{graphicx} % Para insertar imágenes
\usepackage{amsmath, amssymb} % Símbolos matemáticos
\usepackage{hyperref} % Enlaces y referencias clicables
\usepackage{lmodern} % Fuente moderna
\usepackage{geometry} % Márgenes
\usepackage[most]{tcolorbox} % Cajitas para ejemplos y procedimientos
\usepackage{tikz}
\usetikzlibrary{matrix,backgrounds}
\geometry{margin=2.5cm}

% ==== Definición de estilos de caja ====
\tcbset{
  colback=white,
  colframe=black!60,
  sharp corners,
  boxrule=0.8pt,
left=6pt,right=6pt,top=6pt,bottom=6pt
}

\newtcolorbox{procedimiento}[1][]{
  colback=green!5!white,
  colframe=green!60!black,
  title=\textbf{Procedimiento},
  fonttitle=\bfseries,
  #1
}

\newtcolorbox{ejemplo}[1][]{
  colback=blue!5!white,
  colframe=blue!60!black,
  title=\textbf{Ejemplo},
  fonttitle=\bfseries,
  #1
}

% ==== Documento principal ====
\title{Álgebra Lineal y Estructuras Matemáticas (ALEM) Resumen Tema 3}
\author{Alonso Hernández Robles (Grupo E)}
\date{01/12/2025}

\begin{document}

\maketitle

\tableofcontents
\newpage

\section{Matrices}

\begin{procedimiento}[title=Definición de $\mathfrak{M}_{m \times n}(\mathbb{K})$]
$\mathfrak{M}_{m \times n}(\mathbb{K})$ es el conjunto de todas las matrices de orden $m \times n$ con coeficientes en un cuerpo $\mathbb{K}$.
\end{procedimiento}

\begin{procedimiento}[title=Representación matricial]
Una matriz $A = (a_{ij}) \in \mathfrak{M}_{m \times n}(\mathbb{K})$ se representa como:
\[
A =
\begin{pmatrix}
a_{11} & \cdots & a_{1n} \\
\vdots & \ddots & \vdots \\
a_{m1} & \cdots & a_{mn}
\end{pmatrix}.
\]
\end{procedimiento}

\begin{procedimiento}[title=Igualdad de matrices]
Dos matrices $A = (a_{ij})$ y $B = (b_{ij})$ son iguales ($A = B$) si $a_{ij} = b_{ij}$ para todo $i,j$.
\end{procedimiento}

\begin{procedimiento}[title=Submatriz]
Una submatriz de una matriz $A$ es la que resulta de suprimir algunas filas y/o columnas de $A$.
\end{procedimiento}

\begin{ejemplo}[title=Ejemplos de submatrices]
Sea
\[
A=
\begin{pmatrix}
1 & -1 & 4 \\
2 & 3 & \tfrac{1}{2} \\
7 & 0 & \sqrt{2}
\end{pmatrix}
\in \mathfrak{M}_{3\times 3}(\mathbb{R}).
\]

Entonces,
\[
\begin{aligned}
A_1=
\begin{pmatrix}
1 & 4 \\
7 & \sqrt{2}
\end{pmatrix},
\qquad
A_2=
\begin{pmatrix}
1 & 4 \\
2 & \tfrac{1}{2} \\
7 & \sqrt{2}
\end{pmatrix}
\end{aligned}
\]
son submatrices de $A$.
\end{ejemplo}

\begin{procedimiento}[title=Diagonal principal]
La diagonal principal de $A = (a_{ij}) \in \mathfrak{M}_{m \times n}(\mathbb{K})$ es la secuencia $a_{11}, a_{22}, \dots$
\[
\begin{pmatrix}
\boxed{a_{11}} & a_{12} & \cdots \\
a_{21} & \boxed{a_{22}} & \cdots \\
\vdots & \vdots & \ddots &\\
\end{pmatrix}
\]
\end{procedimiento}

\begin{ejemplo}[title=Ejemplos de diagonales principales]
Para
\[
A=
\begin{pmatrix}
\boxed{1} & -1 & 4 \\
2 & \boxed{3} & \tfrac{1}{2} \\
7 & 0 & \boxed{\sqrt{2}}
\end{pmatrix},
\qquad
B=
\begin{pmatrix}
\boxed{1} & 4 \\
2 & \boxed{\tfrac{1}{2}} \\
7 & \sqrt{2}
\end{pmatrix},
\]
la diagonal principal de $A$ es ${1},\, {3},\, {\sqrt{2}}$,  
y la diagonal principal de $B$ es ${1},\, {\tfrac{1}{2}}$.
\end{ejemplo}

\begin{procedimiento}[title=Matriz traspuesta]
La matriz traspuesta de $A \in \mathfrak{M}_{m \times n}(\mathbb{K})$ es la matriz 
$A^T = (b_{ij}) \in \mathfrak{M}_{n \times m}(\mathbb{K})$ tal que 
$b_{ij} = a_{ji}$ para todo $i,j$.  
Se obtiene reflejando $A$ respecto de la diagonal principal.
\end{procedimiento}

\begin{ejemplo}[title=Ejemplo de matriz traspuesta]
Sea
\[
A=
\begin{tikzpicture}[baseline=(m.center)]
    \matrix[matrix of math nodes,left delimiter={(},right delimiter={)}] (m) {
        1 & -1 & 4 \\
        2 & 3 & \tfrac12 \\
        7 & 0 & \sqrt{2} \\
    };

    % Polígono rojo (parte superior de la diagonal)
    \fill[red,opacity=0.125]
        (m-1-1.north west) --
        (m-1-3.north east) --
        (m-3-3.south east) -- cycle;

    % Polígono azul (parte inferior de la diagonal)
    \fill[blue,opacity=0.125]
        (m-1-1.north west) --
        (m-3-1.south west) --
        (m-3-3.south east) -- cycle;
\end{tikzpicture},
\qquad
B=
\begin{tikzpicture}[baseline=(mb.center)]
    \matrix[matrix of math nodes,left delimiter={(},right delimiter={)}] (mb) {
        1 & 4 \\
        2 & \tfrac12 \\
        7 & \sqrt{2} \\
    };

    % Polígono azul (parte superior) – 4 lados
    \fill[red,opacity=0.125]
        (mb-1-1.north west) --
        (mb-1-2.north east) --
        (mb-2-2.center east) -- cycle;

    % Polígono rojo (parte inferior) – 4 lados
    \fill[blue,opacity=0.125]
        (mb-1-1.north west) --
        (mb-3-1.south west) --
        (mb-3-2.south east) --
        (mb-2-2.center east) -- cycle;
\end{tikzpicture}
\]

Entonces,
\[
A^T =
\begin{tikzpicture}[baseline=(mt.center)]
    \matrix[matrix of math nodes,left delimiter={(},right delimiter={)}] (mt) {
        1 & 2 & 7 \\
        -1 & 3 & 0 \\
        4 & \tfrac12 & \sqrt{2} \\
    };

    % Polígono azul (parte superior de la diagonal) — colores invertidos
    \fill[blue,opacity=0.125]
        (mt-1-1.north west) --
        (mt-1-3.north east) --
        (mt-3-3.south east) -- cycle;

    % Polígono rojo (parte inferior de la diagonal)
    \fill[red,opacity=0.125]
        (mt-1-1.north west) --
        (mt-3-1.south west) --
        (mt-3-3.south east) -- cycle;
\end{tikzpicture}
\qquad
B^T =
\begin{tikzpicture}[baseline=(mbt.center)]
    \matrix[matrix of math nodes,left delimiter={(},right delimiter={)}] (mbt) {
        1 & 2 & 7 \\
        4 & \tfrac12 & \sqrt{2} \\
    };

    \fill[blue,opacity=0.125]
        (mbt-1-1.north west) --
        (mbt-1-3.north east) --
        (mbt-2-3.south east) --
        (mbt-2-2.south center) -- cycle;

    % Polígono azul (parte inferior) – 4 lados
    \fill[red,opacity=0.125]
        (mbt-1-1.north west) --
        (mbt-2-2.south center) --
        (mbt-2-1.south west) -- cycle;
    
\end{tikzpicture}
\]
\end{ejemplo}


\begin{procedimiento}[title=Propiedad de la traspuesta]
\[
(A^T)^T = A
\]
\end{procedimiento}

\subsection{Tipos de Matrices}

\begin{procedimiento}[title=Matriz nula]
Una matriz es nula si todos sus coeficientes son cero.
\end{procedimiento}

\begin{ejemplo}[title=Ejemplo de matriz nula]
\[
\begin{pmatrix}
0 & 0 & 0 \\
0 & 0 & 0 \\
0 & 0 & 0
\end{pmatrix}
\in \mathfrak{M}_{3 \times 3}(\mathbb{R}})
\]
es nula.
\end{ejemplo}

\subsubsection{Matrices Cuadradas}

\begin{procedimiento}[title=Matriz cuadrada]
Una matriz $A$ es cuadrada si tiene el mismo número de filas y columnas, es decir, si es de orden $n \times n$.  
En ese caso, es de orden $n$ a secas, y escribimos $A \in \mathfrak{M}_n(\mathbb{K})$.
\end{procedimiento}

\begin{ejemplo}[title=Ejemplo de matriz cuadrada]
\[
A=
\begin{pmatrix}
3 & -2 & 1 \\
0 & 4 & 7 \\
5 & 0 & -1
\end{pmatrix}
\in \mathfrak{M}_3(\mathbb{R})
\]
es cuadrada.
\end{ejemplo}

% -----------------------------------------
\begin{procedimiento}[title=Matriz simétrica]
Una matriz cuadrada $A \in \mathfrak{M}_n(\mathbb{K})$ es simétrica si $A = A^T$.
\end{procedimiento}

\begin{ejemplo}[title=Ejemplo de matriz simétrica]
\[
A=
\begin{pmatrix}
2 & -1 \\
-1 & 7
\end{pmatrix}
=
A^T.
\]
A es simétrica.
\end{ejemplo}

% -----------------------------------------
\begin{procedimiento}[title=Matriz antisimétrica]
Una matriz cuadrada $A \in \mathfrak{M}_n(\mathbb{K})$ es antisimétrica si $a_{ij} = -a_{ji}$ para todo $i,j$.  

Es decir, $A$ es antisimétrica si $A = -A^T$.

\vspace{0.5em}

Como consecuencia de que cuando $i=j$, $a_{ii} = -a_{ii}$:
\begin{itemize}
\item En $\mathbb{R}$, toda diagonal principal debe ser $0$.
\item En $\mathbb{Z}_2$, la diagonal principal debe ser todo $0$ o todo $1$.
\end{itemize}
\end{procedimiento}

\begin{ejemplo}[title=Ejemplo 1 de matriz antisimétrica en $\mathbb{R}$]
\[
A=
\begin{pmatrix}
0 & -2 \\
2 & 0
\end{pmatrix}
\in \mathfrak{M}_2(\mathbb{R})
\]
es antisimétrica.
\end{ejemplo}

\begin{ejemplo}[title=Ejemplo 2 de matriz antisimétrica en $\mathbb{Z}_2$]
\[
B=
\begin{pmatrix}
1 & 0 \\
0 & 1
\end{pmatrix} \in \mathfrak{M}_2(\mathbb{Z}_2), \qquad
C=
\begin{pmatrix}
1 & 1 \\
1 & 1
\end{pmatrix} \in \mathfrak{M}_2(\mathbb{Z}_2).
\]
son antisimétricas.
\end{ejemplo}

% ---------------------------------------
\begin{procedimiento}[title=Matriz triangular superior]
Una matriz cuadrada $A \in \mathfrak{M}_n(\mathbb{K})$ es triangular superior si todo elemento por debajo de la diagonal principal es $0$.  
En ese caso, $A^T$ es triangular inferior.
\end{procedimiento}

\begin{ejemplo}[title=Ejemplo de matriz triangular superior]
\[
A=
\begin{pmatrix}
2 & 3 & -2 \\
0 & 1 & 1 \\
0 & 0 & 4
\end{pmatrix}
\]
es triangular superior.
\end{ejemplo}



% -----------------------------------------
\begin{procedimiento}[title=Matriz triangular inferior]
Una matriz cuadrada $A \in \mathfrak{M}_n(\mathbb{K})$ es triangular inferior si todo elemento por encima de la diagonal principal es $0$.  
En ese caso, $A^T$ es triangular superior.
\end{procedimiento}

\begin{ejemplo}[title=Ejemplo de matriz triangular inferior]
\[
A=
\begin{pmatrix}
-2 & 0 & 0 \\
1 & 7 & 0 \\
8 & 0 & 5
\end{pmatrix}
\]
es triangular inferior.
\end{ejemplo}



% -----------------------------------------
\begin{procedimiento}[title=Matriz diagonal]
Una matriz cuadrada $A \in \mathfrak{M}_n(\mathbb{K})$ es diagonal si es triangular superior e inferior simultáneamente.  
Es decir, todo elemento fuera de la diagonal principal es $0$.
\end{procedimiento}

\begin{ejemplo}[title=Ejemplos de matrices diagonales]
\[
A=
\begin{pmatrix}
2 & 0 & 0 \\
0 & -1 & 0 \\
0 & 0 & 7
\end{pmatrix},
\quad
B=
\begin{pmatrix}
1 & 0 & 0 \\
0 & 0 & 0 \\
0 & 0 & -1
\end{pmatrix},
\quad
C=
\begin{pmatrix}
0 & 0 & 0 \\
0 & 0 & 0 \\
0 & 0 & 0
\end{pmatrix}
\]
son matrices diagonales.
\end{ejemplo}

\begin{procedimiento}[title=Matriz identidad]
La matriz identidad de orden $n$ es la matriz diagonal $I_n$ donde todos los elementos de la diagonal principal son 1, es decir:
\[
I_n =
\begin{pmatrix}
1 & 0 & \cdots & 0 \\
0 & 1 & \cdots & 0 \\
\vdots & \vdots & \ddots & \vdots \\
0 & 0 & \cdots & 1
\end{pmatrix} \in \mathfrak{M}_n(\mathbb{K}).
\]
\end{procedimiento}

\section{Operaciones con Matrices}

\subsection{Suma de matrices}

\begin{procedimiento}[title=Suma]
Sean $A = (a_{ij})$, $B = (b_{ij}) \in \mathfrak{M}_{m \times n}(\mathbb{K})$. La suma de $A$ y $B$ es la matriz $A+B = (c_{ij}) \in \mathfrak{M}_{m \times n}(\mathbb{K})$ donde $c_{ij} = a_{ij} + b_{ij}$ para todo $i,j$.
\end{procedimiento}

\begin{ejemplo}[title=Ejemplo de suma de matrices]
Sea
\[
A=
\begin{pmatrix}
5 & 2 & 4 \\
1 & 3 & 6
\end{pmatrix},
\quad
B=
\begin{pmatrix}
6 & 5 & 1 \\
3 & 2 & 0
\end{pmatrix}
\in \mathfrak{M}_{2 \times 3}(\mathbb{Z}_7).
\]

Entonces
\[
A + B =
\begin{pmatrix}
4 & 0 & 5 \\
4 & 5 & 6
\end{pmatrix}.
\]
\end{ejemplo}

\begin{procedimiento}[title=Matriz opuesta]
La matriz opuesta de $A \in \mathfrak{M}_{m \times n}(\mathbb{K})$ es la matriz $-A \in \mathfrak{M}_{m \times n}(\mathbb{K})$ que cumple
\[
A + (-A) =
\begin{pmatrix}
0 & \cdots & 0 \\
\vdots & \ddots & \vdots \\
0 & \cdots & 0
\end{pmatrix} \in \mathfrak{M}_{m \times n}(\mathbb{K}) \quad \text{(matriz nula)}.
\]
\end{procedimiento}

\begin{procedimiento}[title=Propiedades de la suma]
La suma verifica:
\begin{itemize}
\item Conmutativa: $A + B = B + A$.
\item Asociativa: $(A + B) + C = A + (B + C)$.
\item Elemento neutro: la matriz nula de orden $m \times n$.
\end{itemize}
El conjunto $\big(\mathfrak{M}_{m \times n}(\mathbb{K}), +\big)$ es un grupo abeliano o conmutativo.
\begin{itemize}
    \item $(A+B)^T = A^T + B^T$
\end{itemize}
\end{procedimiento}


\subsection{Producto de un escalar}

\begin{procedimiento}[title=Producto de un escalar]
Sea $\lambda \in \mathbb{K}$ y $A = (a_{ij}) \in \mathfrak{M}_{m \times n}(\mathbb{K})$.  
El producto $\lambda A = (b_{ij}) \in \mathfrak{M}_{m \times n}(\mathbb{K})$ se define por $b_{ij} = \lambda a_{ij}$ para todo $i,j$.
\end{procedimiento}

\begin{ejemplo}[title=Ejemplo de producto por un escalar]
Sea
\[
A=
\begin{pmatrix}
5 & 2 & 4 \\
1 & 3 & 6
\end{pmatrix} \in \mathfrak{M}_{2\times3}(\mathbb{Z}_7)
\]

Entonces
\[
3A =
\begin{pmatrix}
1 & 6 & 5 \\
3 & 2 & 4
\end{pmatrix}, 
\quad
0A =
\begin{pmatrix}
0 & 0 & 0 \\
0 & 0 & 0
\end{pmatrix}.
\]
\end{ejemplo}

\begin{procedimiento}[title=Matriz opuesta]
$(-1)A = -A$ da la matriz opuesta.
\end{procedimiento}

\begin{procedimiento}[title=Propiedades del producto de un escalar]
Para todo $A \in \mathfrak{M}_{m \times n}(\mathbb{K})$ y para todo $\lambda \in \mathbb{K}$:
\begin{itemize}
\item $(\lambda A)^T = \lambda A^T$
\end{itemize}
\end{procedimiento}

\subsection{Producto de matrices}

\begin{procedimiento}[title=Producto de matrices]
Sea $A = (a_{ik}) \in \mathfrak{M}_{m \times t}(\mathbb{K})$ y $B = (b_{kj}) \in \mathfrak{M}_{t \times n}(\mathbb{K})$.  
El producto de las matrices $A$ y $B$ es
\[
A \cdot B = (c_{ij}) \in \mathfrak{M}_{m \times n}(\mathbb{K}),
\quad \text{donde} \quad
c_{ij} = \sum_{k=1}^{t} a_{ik} b_{kj}.
\]

\[
A \cdot B =
\begin{pmatrix}
a_{i1} & \cdots & a_{it}
\end{pmatrix}
\begin{pmatrix}
b_{1j} \\ \vdots \\ b_{tj}
\end{pmatrix}
=
(c_{ij}).
\]
\end{procedimiento}

\begin{ejemplo}[title=Ejemplo de producto de matrices]
Sea
\[
A =
\begin{pmatrix}
2 & 0 \\
4 & 3
\end{pmatrix} \in \mathfrak{M}_{2 \times 2}(\mathbb{Z}_5), 
\quad
B =
\begin{pmatrix}
1 & 2 & 3 \\
2 & 4 & 1
\end{pmatrix} \in \mathfrak{M}_{2 \times 3}(\mathbb{Z}_5).
\]

El producto $A \cdot B$ se calcula dentro de $\mathbb{Z}_5$ como:
\[
A \cdot B =
\begin{pmatrix}
2\cdot 1 + 0 \cdot 2 & 2\cdot 2 + 0 \cdot 4 & 2\cdot 3 + 0 \cdot 1 \\
4\cdot 1 + 3 \cdot 2 & 4\cdot 2 + 3 \cdot 4 & 4\cdot 3 + 3 \cdot 1
\end{pmatrix} =
\begin{pmatrix}
2 & 4 & 1 \\
0 & 0 & 0
\end{pmatrix} \in \mathfrak{M}_{2 \times 3}(\mathbb{Z}_5).
\]

$B \cdot A$ no está definido porque nº columnas de B $= 3 \neq 2 = $ nº filas de A.
\end{ejemplo}

\begin{procedimiento}[title=Propiedades del producto de matrices]
\begin{itemize}
\item Para todo $A \in \mathfrak{M}_{m \times n}(\mathbb{K})$, $I_m \cdot A = A = A \cdot I_n$.
\item Asociativa: $A \cdot (B \cdot C) = (A \cdot B) \cdot C$.
\item Distributiva: 
\[
A \cdot (B_1 + B_2) = A \cdot B_1 + A \cdot B_2, \quad 
(B_1 + B_2) \cdot C = B_1 \cdot C + B_2 \cdot C.
\]
\item Para todo $\lambda \in \mathbb{K}$, $(\lambda A) \cdot B = \lambda (A \cdot B) = A \cdot (\lambda B)$.
\item En general, el producto de matrices no es conmutativo: $A \cdot B \neq B \cdot A$.

Si $A \cdot B = B \cdot A$, se cumplen las identidades notables:
\begin{itemize}
\item $(A + B)^2 = A^2 + 2AB + B^2$
\item $(A - B)^2 = A^2 - 2AB + B^2$
\item $(A + B)(A - B) = A^2 - B^2$
\end{itemize}
\item Traspuesta: $(A \cdot B)^T = B^T \cdot A^T$.
\end{itemize}
\end{procedimiento}

\section{Operaciones o Transformaciones Elementales sobre Matrices}

\subsection{Operaciones por filas}

\begin{procedimiento}[title=Fila nula]
Una fila en una matriz es nula si está formada íntegramente por ceros.
\end{procedimiento}

\begin{ejemplo}[title=Ejemplo de fila nula]
\[
\begin{pmatrix}
0 & 0 & 0 & 0 \\
1 & 2 & 3 & 4
\end{pmatrix}
\]
La primera fila es nula, la segunda no lo es.
\end{ejemplo}

\begin{procedimiento}[title=Pivote o término líder]
En una fila no nula, el pivote o término líder de la fila es el primer elemento distinto de cero que aparezca por la izquierda.
\end{procedimiento}

\begin{ejemplo}[title=Ejemplo de pivote]
\[
\begin{pmatrix}
0 & \boxed{2} & 3 & 4 \\
0 & 0 & 0 & \boxed{1} \\
0 & 0 & 0 & 0
\end{pmatrix}
\]

\begin{itemize}
\item 2 es el pivote de la primera fila.
\item 1 es el pivote de la segunda fila.
\end{itemize}
\end{ejemplo}

\begin{procedimiento}[title=Forma escalonada por filas]
Una matriz $A \in \mathfrak{M}_{m \times n}(\mathbb{K})$ está en forma escalonada por filas si verifica:
\begin{enumerate}
\item Las filas nulas, si hay, están todas en la parte inferior.
\item El pivote de una fila está situado más a la derecha que el pivote de cualquier fila superior.
\end{enumerate}
\end{procedimiento}

\begin{ejemplo}[title=Ejemplos de forma escalonada por filas]
\[
\begin{pmatrix}
0 & 2 & 3 & 4 \\
0 & 0 & 0 & 1 \\
0 & 0 & 0 & 0
\end{pmatrix}
\quad \text{está en forma escalonada por filas.}
\]

\[
\begin{pmatrix}
0 & 2 & 3 & 4 \\
0 & 0 & 0 & 0 \\
0 & 0 & 0 & 1
\end{pmatrix}
\quad
\text{\parbox[t]{1\textwidth}{no está en forma escalonada por filas porque la \\ única fila nula no está en la parte inferior.}}
\]

\[
\begin{pmatrix}
2 & 3 & 4 \\
0 & \boxed{2} & 5 \\
0 & \boxed{1} & 3
\end{pmatrix}
\quad
\text{\parbox[t]{1\textwidth}{no está en forma escalonada por filas porque el pivote 1 \\ no está más a la derecha que el 2, pivote de la fila superior.}}
\]
\end{ejemplo}

\begin{procedimiento}[title=Forma escalonada reducida por filas]
Una matriz $A \in \mathfrak{M}_{m \times n}(\mathbb{K})$ está en forma escalonada reducida por filas si cumple:
\begin{enumerate}
\item $A$ está en forma escalonada por filas.
\item Para cada pivote, todos los elementos por encima y por debajo son cero.
\item Todos los pivotes son 1.
\end{enumerate}

Se cumple que:
\begin{itemize}
\item Cualquier matriz nula está en forma escalonada reducida por filas.
\item La matriz identidad $I_n$ está en forma escalonada reducida por filas.
\end{itemize}
\end{procedimiento}

\begin{ejemplo}[title=Ejemplos de forma escalonada reducida por filas]

\[
\begin{pmatrix}
0 & 2 & 3 & 0 \\
0 & 0 & 0 & 1 \\
0 & 0 & 0 & 0
\end{pmatrix}
\quad
\text{\parbox[t]{1\textwidth}{no está en forma escalonada reducida por filas, \\ porque los pivotes no son todos 1.}}
\]

\[
\begin{pmatrix}
0 & 1 & 3 & 4 \\
0 & 0 & 0 & 1 \\
0 & 0 & 0 & 0
\end{pmatrix}
\quad
\text{\parbox[t]{1\textwidth}{no está en forma escalonada reducida por filas, \\ porque el pivote de la segunda fila no tiene \\ todos los elementos por encima y por debajo cero.}}
\]

\[
\begin{pmatrix}
0 & \boxed{1} & 4 & 0 \\
0 & 0 & 0 & \boxed{1} \\
0 & 0 & 0 & 0
\end{pmatrix}
\quad
\text{\parbox[t]{1\textwidth}{está en forma escalonada reducida por filas.}}
\]

\end{ejemplo}

\begin{procedimiento}[title=Operación o Transformación de filas]
Para $A \in \mathfrak{M}_{m \times n}(\mathbb{K})$, una operación o transformación elemental de filas puede ser:
\begin{enumerate}
\item Intercambiar de posición dos filas:
\[
F_i \longleftrightarrow F_j
\]
\item Multiplicar una fila por un escalar distinto de cero:
\[
F_i \longleftarrow \lambda F_j \qquad (\lambda F_j)
\]
\item Sumarle a una fila otra multiplicada por un escalar:
\[
F_i \longleftarrow F_i + \lambda F_j \qquad (F_i + \lambda F_j)
\]
\end{enumerate}
\end{procedimiento}

\begin{procedimiento}[title=Matrices equivalentes por filas]
Dos matrices $A,B \in \mathfrak{M}_{m \times n}(\mathbb{K})$ son equivalentes por filas, denotado
\[
A \sim_f B,
\]
si $A=B$ o si es posible obtener $B$ a partir de $A$ aplicando un número finito de operaciones elementales de filas. También lo son si $H_f(A) = H_f(B)$. 
\end{procedimiento}

\begin{procedimiento}[title=Forma normal de Hermite por filas]
Toda matriz $A \in \mathfrak{M}_{m \times n}(\mathbb{K})$ es equivalente por filas a una única matriz en forma escalonada reducida por filas
\[
H_f(A),
\]
llamada la \textbf{Forma normal de Hermite por filas} de A.
\end{procedimiento}

\begin{procedimiento}[title=Rango por filas]
El rango por filas de $A$, denotado $\operatorname{Rg}(A)$, es el número de filas no nulas en $H_f(A)$.

Equivalentemente, es también el número de filas no nulas de cualquier matriz $B$ en forma escalonada por filas tal que $A \sim_f B$.
\end{procedimiento}

\begin{ejemplo}[title=Ejemplo de cálculo del rango por filas]
Sea
\[
A =
\begin{pmatrix}
6 & 2 & 1 & 4 \\
5 & 3 & 1 & 3 \\
4 & 6 & 3 & 5
\end{pmatrix} \in \mathfrak{M}_{3 \times 4}(\mathbb{Z}_7).
\]

Aplicamos operaciones elementales de fila para obtener una matriz en forma escalonada por filas:

\[
A \overset{\substack{F_2 + 5 F_1 \\ F_3 + 4 F_1\\{}}}{\sim_f}
\begin{pmatrix}
6 & 2 & 1 & 4 \\
0 & 6 & 6 & 2 \\
0 & 0 & 0 & 0
\end{pmatrix}
\]

Esta matriz está en forma escalonada por filas, por lo que
\[
\operatorname{Rg}(A) = \text{nº filas no nulas} = 2.
\]

Ahora hallamos $H_f(A)$, la forma escalonada reducida por filas, mediante más operaciones:

\[
\begin{pmatrix}
6 & 2 & 1 & 4 \\
0 & 6 & 6 & 2 \\
0 & 0 & 0 & 0
\end{pmatrix}
\overset{\substack{F_1 +2F_2 \\{}}}{\sim_f}
\begin{pmatrix}
6 & 0 & 6 & 1 \\
0 & 6 & 6 & 2 \\
0 & 0 & 0 & 0
\end{pmatrix}
\overset{\substack{-F_1\\-F_2\\{}}}{\sim_f}
\begin{pmatrix}
1 & 0 & 1 & 6 \\
0 & 1 & 1 & 5 \\
0 & 0 & 0 & 0
\end{pmatrix}
= H_f(A)
\]

Así obtenemos la forma normal de Hermite por filas de $A$ y comprobamos que $\operatorname{Rg}(A) = 2$.
\end{ejemplo}

\subsection{Operaciones por columnas}

\begin{procedimiento}[title=Columna nula]
Una columna en una matriz es nula si está formada íntegramente por ceros.
\end{procedimiento}

\begin{procedimiento}[title=Pivote o término líder]
En una columna no nula, el pivote o término líder de la columna es el primer elemento distinto de cero que aparezca por arriba.
\end{procedimiento}

\begin{procedimiento}[title=Forma escalonada por columnas]
Una matriz $A \in \mathfrak{M}_{m \times n}(\mathbb{K})$ está en forma escalonada por columnas si verifica:
\begin{enumerate}
\item Las columnas nulas, si hay, están todas en la parte derecha.
\item El pivote de una columna está situado más abajo que el pivote de cualquier columna anterior.
\end{enumerate}
\end{procedimiento}

\begin{procedimiento}[title=Forma escalonada reducida por columnas]
Una matriz $A \in \mathfrak{M}_{m \times n}(\mathbb{K})$ está en forma escalonada reducida por columnas si cumple:
\begin{enumerate}
\item $A$ está en forma escalonada por columnas.
\item Para cada pivote, todos los elementos a la izquierda y a la derecha son cero.
\item Todos los pivotes son 1.
\end{enumerate}

Se cumple que:
\begin{itemize}
\item Cualquier matriz nula está en forma escalonada reducida por columnas.
\item La matriz identidad $I_n$ está en forma escalonada reducida por columnas.
\end{itemize}
\end{procedimiento}

\begin{procedimiento}[title=Operación o Transformación de columnas]
Para $A \in \mathfrak{M}_{m \times n}(\mathbb{K})$, una operación o transformación elemental de columnas puede ser:
\begin{enumerate}
\item Intercambiar de posición dos columnas:
\[
C_i \longleftrightarrow C_j
\]
\item Multiplicar una columna por un escalar distinto de cero:
\[
C_i \longleftarrow \lambda C_j \qquad (\lambda C_j)
\]
\item Sumarle a una columna otra multiplicada por un escalar:
\[
C_i \longleftarrow C_i + \lambda C_j \qquad (C_i + \lambda C_j)
\]
\end{enumerate}
\end{procedimiento}

\begin{procedimiento}[title=Matrices equivalentes por columnas]
Dos matrices $A,B \in \mathfrak{M}_{m \times n}(\mathbb{K})$ son equivalentes por columnas, denotado
\[
A \sim_c B,
\]
si $A=B$ o si es posible obtener $B$ a partir de $A$ aplicando un número finito de operaciones elementales de columnas. También lo son si $H_c(A) = H_c(B)$.
\end{procedimiento}

\begin{procedimiento}[title=Forma normal de Hermite por columnas]
Toda matriz $A \in \mathfrak{M}_{m \times n}(\mathbb{K})$ es equivalente por columnas a una única matriz en forma escalonada reducida por columnas
\[
H_c(A),
\]
llamada la \textbf{Forma normal de Hermite por columnas} de A. 
\end{procedimiento}

\begin{procedimiento}[title=Rango por columnas]
El rango por columnas de $A$, denotado $\operatorname{Rg}(A)$, es el número de columnas no nulas en $H_c(A)$.

Equivalentemente, es también el número de columnas no nulas de cualquier matriz $B$ en forma escalonada por columnas tal que $A \sim_c B$.
\end{procedimiento}

\begin{ejemplo}[title=Ejemplo de cálculo del rango por filas]
Sea
\[
A =
\begin{pmatrix}
1 & 4 & 1 \\
1 & 1 & 2 \\
6 & 5 & 3 \\
1 & 1 & 1
\end{pmatrix}
\in \mathfrak{M}_{4 \times 3}(\mathbb{Z}_7).
\]

Aplicamos operaciones elementales de columna para obtener una matriz en forma escalonada por columnas:

\[
A
\overset{\substack{
C_2 + 3C_1 \\
C_3 - C_1 \\ {}
}}{\sim_c}
\begin{pmatrix}
1 & 0 & 0 \\
1 & 4 & 1 \\
6 & 2 & 4 \\
1 & 4 & 0
\end{pmatrix}
\overset{\substack{
2C_2 \\ {}
}}{\sim_c}
\begin{pmatrix}
1 & 0 & 0 \\
1 & 1 & 1 \\
6 & 4 & 4 \\
1 & 1 & 0
\end{pmatrix}
\overset{\substack{
C_3 - C_2 \\ {}
}}{\sim_c}
\begin{pmatrix}
1 & 0 & 0 \\
1 & 1 & 0 \\
6 & 4 & 0 \\
1 & 1 & 6
\end{pmatrix}
\]

Esta matriz está en forma escalonada por columnas, por lo que:
\[
\operatorname{Rg}(A) = \text{nº columnas no nulas} = 3.
\]

Ahora hallamos $H_c(A)$, la forma escalonada reducida por columnas:

\[
\begin{pmatrix}
1 & 0 & 0 \\
1 & 1 & 0 \\
6 & 4 & 0 \\
1 & 1 & 6
\end{pmatrix}
\overset{\substack{
C_1 - C_2 \\
- C_3 \\ {}
}}{\sim_c}
\begin{pmatrix}
1 & 0 & 0 \\
0 & 1 & 0 \\
2 & 4 & 0 \\
0 & 1 & 1
\end{pmatrix}
\overset{\substack{
C_2 - C_3 \\ {}
}}{\sim_c}
\begin{pmatrix}
1 & 0 & 0 \\
0 & 1 & 0 \\
2 & 4 & 0 \\
0 & 0 & 1
\end{pmatrix}
= H_c(A)
\]

Así obtenemos la forma normal de Hermite por columnas de $A$ y comprobamos que $\operatorname{Rg}(A)=3$.
\end{ejemplo}

\subsection{Teorema del Rango}

\begin{procedimiento}[title=Teorema del rango]
Para toda $A \in \mathfrak{M}_{m \times n}(\mathbb{K})$, el rango por filas de $A$ y el
rango por columnas de $A$ son el mismo, y ambos son $\operatorname{Rg}(A)$.
\end{procedimiento}

\section{Sistemas de Ecuaciones Lineales}

\begin{procedimiento}[title=Sistema de ecuaciones lineales]
Un sistema de $m$ ecuaciones con $n$ incógnitas, donde $m,n \in \mathbb{N}$, ambos $\ge 1$, es de la forma:
\[
\begin{cases}
a_{11}x_1 + a_{12}x_2 + \dots + a_{1n}x_n = b_1 \\
\,\,\,\,\,\,\,\vdots\,\,\,\,\,\,\,\,\,\,\,\,\,\,\,\,\,\,\,\,\vdots\,\,\,\,\,\,\,\,\,\,\,\,\ddots\,\,\,\,\,\,\,\,\,\,\,\,\,\vdots\,\,\,\,\,\,\,\,\,\,\,\,\,\,\,\,\vdots \\
a_{m1}x_1 + a_{m2}x_2 + \dots + a_{mn}x_n = b_m
\end{cases}
\]
donde $a_{ij}, b_i \in \mathbb{K}$.

\begin{itemize}
    \item Los $a_{ij}$ son los coeficientes del sistema.
    \item Los $b_i$ son los términos independientes.
    \item $x_1,x_2,\dots,x_n$ son las incógnitas del sistema.
\end{itemize}

La solución del sistema es una $n$-tupla $(\alpha_1,\alpha_2,\dots,\alpha_n) \in \mathbb{K}^n$ tal que, al sustituir $x_i=\alpha_i$, se satisfacen todas las ecuaciones.
\end{procedimiento}

\begin{procedimiento}[title=Matrices del sistema]
La matriz de coeficientes es
\[
A =
\begin{pmatrix}
a_{11} & a_{12} & \dots & a_{1n}\\
\vdots & \vdots & \ddots & \vdots\\
a_{m1} & a_{m2} & \dots & a_{mn}
\end{pmatrix}
\in \mathfrak{M}_{m \times n}(\mathbb{K}).
\]

La matriz de términos independientes es
\[
B =
\begin{pmatrix}
b_1 \\ \vdots \\ b_m
\end{pmatrix}
\in \mathfrak{M}_{m \times 1}(\mathbb{K}).
\]

La matriz de incógnitas es
\[
X =
\begin{pmatrix}
x_1 \\ \vdots \\ x_n
\end{pmatrix}.
\]

La expresión matricial del sistema es simplemente:
\[
A\,X = B.
\]
\end{procedimiento}

\begin{ejemplo}[title=Ejemplo de sistema y su expresión matricial]
Consideremos el sistema:
\[
\begin{cases}
3x_1 - 4x_2 + 8x_3 = 5 \\
2x_1 \;\;\;\;\;\;\;\;\; - \tfrac{1}{2}x_3 = 0
\end{cases}
\]

Entonces:
\[
A =
\begin{pmatrix}
3 & -4 & 8 \\
2 & 0 & -\tfrac12
\end{pmatrix},
\qquad
B =
\begin{pmatrix}
5 \\ 0
\end{pmatrix},
\qquad
X =
\begin{pmatrix}
x_1 \\ x_2 \\ x_3
\end{pmatrix}.
\]

Por tanto, el sistema se escribe como $A\,X = B$.
\end{ejemplo}

\begin{procedimiento}[title=Matriz ampliada del sistema]
La matriz ampliada del sistema $A\,X=B$ es la matriz
\[
(A \mid B)
=
\begin{pmatrix}
a_{11} & a_{12} & \dots & a_{1n} & \Big| & b_1 \\
\vdots & \vdots & \ddots & \vdots & \Big| & \vdots \\
a_{m1} & a_{m2} & \dots & a_{mn} & \Big| & b_m
\end{pmatrix},
\]
que reúne en una sola matriz la información de los coeficientes y los términos independientes.
\end{procedimiento}

\begin{ejemplo}[title=Ejemplo de matriz ampliada]
Para el sistema
\[
\begin{cases}
3x_1 - 4x_2 + 8x_3 = 5 \\
2x_1 \;\;\;\;\;\;\;\;\; - \tfrac12 x_3 = 0
\end{cases}
\]
tenemos:

\[
A=
\begin{pmatrix}
3 & -4 & 8 \\
2 & 0 & -\tfrac12
\end{pmatrix},
\qquad
B =
\begin{pmatrix}
5 \\ 0
\end{pmatrix}.
\]

Entonces la matriz ampliada es:
\[
(A \mid B) =
\begin{pmatrix}
3 & -4 & 8 & \big| & 5 \\
2 & 0 & -\tfrac12 & \big| & 0
\end{pmatrix}.
\]
\end{ejemplo}

\begin{procedimiento}[title=Tipos de sistemas]
\begin{itemize}
    \item Sistema Incompatible (SI): no tiene solución.
    \item Sistema Compatible
    \begin{itemize}
        \item Sistema Compatible Determinado (SCD): 1 solución.
        \item Sistema Compatible Indeterminado (SCI): más de 1 solución.
    \end{itemize}
\end{itemize}
\end{procedimiento}

\begin{procedimiento}[title=Sistemas equivalentes]
Dos sistemas de ecuaciones lineales sobre las mismas incógnitas son equivalentes si tienen exactamente las mismas soluciones.

\vspace{0.5em}

Sean $S$ y $S'$ sistemas de ecuaciones lineales cuyas matrices ampliadas son, 
respectivamente, $(A\,|\,B)$ y $(A'\,|\,B')$.  
$S$ y $S'$ son equivalentes si cumplen:
\[
(A\,|\,B)\ \sim_f\ (A'\,|\,B'),
\]
\end{procedimiento}

\begin{procedimiento}[title=Resolución de sistemas por Métodos de Gauss y Gauss-Jordan]

Para resolver un sistema $S$, procedemos así:
\begin{itemize}
    \item Partimos de la matriz ampliada $(A\,|\,B)$.
    \item Aplicamos operaciones elementales de filas hasta obtener una matriz ampliada equivalente $(A'\,|\,B')$ (de un sistema equivalente $S'$) que esté en \textbf{forma escalonada por filas}.
    \item Una vez obtenida la forma escalonada por filas, tenemos dos opciones:
    \begin{itemize}
        \item \textbf{Método de Gauss}: Aplicar el \textbf{Teorema de Rouché--Frobenius}.
        \item \textbf{Método de Gauss-Jordan}: Seguir aplicando operaciones elementales de filas a $(A'\,|\,B')$ hasta obtener la forma normal de Hermite por filas de $(A'\,|\,B')$, es decir, una matriz en forma escalonada \textbf{reducida} por filas.
    \end{itemize}
    \item En ambos casos, averiguar el valor de todas las incógnitas como se explica en la siguiente subsección.
\end{itemize}
\end{procedimiento}

\subsection{Teorema de Rouché--Frobenius}

\begin{procedimiento}[title=Teorema de Rouché--Frobenius]
    
\begin{itemize}
    \item[1)] El sistema $A X = B$ es \textbf{incompatible} si y sólo si aparece una ecuación del tipo  
    \[
    0x_1 + 0x_2 + \cdots + 0x_n = b_i,\quad b_i \neq 0.
    \]

    \item[2)] El sistema $A X = B$ es \textbf{compatible} si y sólo si 
    \[
    \operatorname{Rg}(A) = \operatorname{Rg}(A\,|\,B).
    \]

    \item[3)] El sistema $A X = B$ es \textbf{compatible determinado} si y sólo si 
    \[
    \operatorname{Rg}(A) = \operatorname{Rg}(A\,|\,B) = \text{nº de incógnitas}.
    \]
\end{itemize}

\end{procedimiento}

\begin{procedimiento}[title=Cómo resolver un sistema compatible]
Una vez obtenida la matriz ampliada $(A'\,|\,B')$ en forma escalonada por filas, procedemos así:

\begin{enumerate}
    \item Identificar las columnas con pivote: sus incógnitas son las \textbf{principales}.
    \item El resto de incógnitas son \textbf{secundarias} o \textbf{parámetros}.
    \item Reescribir cada ecuación dejando las incógnitas principales a la izquierda y las secundarias a la derecha.
    \item Elegir valores libres para los parámetros: cada elección determina un único valor para cada incógnita principal.
    \item Conclusión:
    \begin{itemize}
        \item Si aparecen parámetros, el sistema es SCI.
        \item Si todas las incógnitas son principales (no hay parámetros), el sistema es SCD.
    \end{itemize}
\end{enumerate}
\end{procedimiento}

\begin{ejemplo}[title=Resolución de un sistema lineal sobre $\mathbb{Z}_7$]
Resolver el sistema:
\[
\begin{cases}
2x + 5y + 4z = 4 \\
5x + 6y + 2z = 3 \\
3x + 2y + 3z = 6 
\end{cases}
\qquad \text{sobre } \mathbb{Z}_7.
\]

La matriz ampliada es:
\[
(A\,|\,B)=
\begin{pmatrix}
2 & 5 & 4 \;\Big|\; 4 \\
5 & 6 & 2 \;\Big|\; 3 \\
3 & 2 & 3 \;\Big|\; 6
\end{pmatrix}.
\]

Aplicamos operaciones elementales de fila:

\[
(A\,|\,B)
\overset{\substack{F_2 + F_1 \\ F_3 - 3F_1 \\{}}}{\sim_f}
\begin{pmatrix}
2 & 5 & 4 \;\Big|\; 4 \\
0 & 4 & 6 \;\Big|\; 0 \\
0 & 5 & 4 \;\Big|\; 0
\end{pmatrix}
\overset{\substack{2F_2 \\{}}}{\sim_f}
\begin{pmatrix}
2 & 5 & 4 \;\Big|\; 4 \\
0 & 1 & 5 \;\Big|\; 0 \\
0 & 5 & 4 \;\Big|\; 0
\end{pmatrix}
\overset{\substack{F_3 + 2F_2 \\{}}}{\sim_f}
\begin{pmatrix}
2 & 5 & 4 \;\Big|\; 4 \\
0 & 1 & 5 \;\Big|\; 0 \\
0 & 0 & 0 \;\Big|\; 0
\end{pmatrix}.
\]

\vspace{0.5em}

Usamos el método de \textbf{Gauss}. Observando el nº filas no nulas en $A$ y $A\,|\,B$:

\[
\operatorname{Rg}(A)=\operatorname{Rg}(A\,|\,B)=2 \neq 3 = \text{nº incógnitas}.
\]

Por tanto, el sistema es \textbf{SCI}.

\bigskip

El sistema equivalente queda:
\[
\begin{cases}
2x + 5y + 4z = 4,\\[2mm]
\phantom{2x +{}} y + 5z = 0.
\end{cases}
\]

Las incógnitas principales son \(x,y\).

La incógnita secundaria (parámetro) es \(z\).

\[
\begin{cases}
2x + 5y = 4 - 4z \\
y = -5 z \\
z = \lambda
\end{cases}
\]

Sustituyendo:

\[
\begin{cases}
y = 2 \lambda \\
2x = 4 -5y -4z \to 2x = 4 + 2(2\lambda) + 3 \lambda \to 2x = 4 + 0 \lambda \to x = 2
\end{cases}
\]

\[
\boxed{
x = 2,\qquad y = 2\lambda,\qquad z = \lambda,\qquad \lambda \in \mathbb{Z}_7.
}
\]

\end{ejemplo}

\begin{ejemplo}[title=Resolución de un sistema lineal sobre $\mathbb{Z}_7$]
Otro camino es usar el método de \textbf{Gauss-Jordan}, es decir, obteniendo $H_f(A\,|\,B)$:

\[
\begin{pmatrix}
2 & 5 & 4 \;\Big|\; 4 \\
0 & 1 & 5 \;\Big|\; 0 \\
0 & 0 & 0 \;\Big|\; 0
\end{pmatrix}
\overset{\substack{F_1 + 2 F_2 \\{}}}{\sim_f}
\begin{pmatrix}
2 & 0 & 0 \;\Big|\; 4 \\
0 & 1 & 5 \;\Big|\; 0 \\
0 & 0 & 0 \;\Big|\; 0
\end{pmatrix}
\overset{\substack{4 F_1 \\{}}}{\sim_f}
\begin{pmatrix}
1 & 0 & 0 \;\Big|\; 2 \\
0 & 1 & 5 \;\Big|\; 0 \\
0 & 0 & 0 \;\Big|\; 0
\end{pmatrix}
= H_f(A\,|\,B).
\]

\[
\Rightarrow
\begin{cases}
x = 2 \\
y + 5 z = 0
\end{cases}
\]

\[
\Rightarrow y = -5 z
\]

\(z = \lambda\). Entonces:

\[
\boxed{
x = 2,\qquad y = 2 \lambda,\qquad z = \lambda,\qquad \lambda \in \mathbb{Z}_7.
}
\]

\end{ejemplo}

\begin{ejemplo}[title=Ejemplo de sistema dependiente de un parámetro]
Resolver el sistema sobre $\mathbb{Z}_7$:

\[
\begin{cases}
3x + 5y + 2z = 6 \\
5x + 6y + z = a \\
2x + y + 6z = 4
\end{cases}
\]

Matriz ampliada:
\[
(A\,|\,B) =
\begin{pmatrix}
3 & 5 & 2 \;\Big|\; 6 \\
5 & 6 & 1 \;\Big|\; a \\
2 & 1 & 6 \;\Big|\; 4
\end{pmatrix}.
\]

Aplicando operaciones elementales de fila:

\[
(A\,|\,B) =
\begin{pmatrix}
3 & 5 & 2 \;\Big|\; 6 \\
5 & 6 & 1 \;\Big|\; a \\
2 & 1 & 6 \;\Big|\; 4
\end{pmatrix}
\overset{\substack{5 F_1 \\{}}}{\sim_f}
\begin{pmatrix}
1 & 4 & 3 \;\Big|\; 2 \\
5 & 6 & 1 \;\Big|\; a \\
2 & 1 & 6 \;\Big|\; 4
\end{pmatrix}
\overset{\substack{F_2 + 2 F_1 \\ F_3 - 2 F_1 \\{}}}{\sim_f}
\begin{pmatrix}
1 & 4 & 3 &\Big| &2 \\
0 & 0 & 0 &\Big| &a+4 \\
0 & 0 & 0 &\Big| &0
\end{pmatrix}.
\]

\vspace{0.5em}

Usamos el método de Gauss:

\[
\operatorname{Rg}(A) = 1, \quad
\operatorname{Rg}(A\,|\,B) =
\begin{cases}
1 & \text{si } a+4 = 0 \iff a = 3,\\
2 & \text{si } a+4 \neq 0 \iff a \neq 3.
\end{cases}
\]

\begin{itemize}
    \item Si $a \neq 3$: $\operatorname{Rg}(A) = 1 \neq 2 = \operatorname{Rg}(A\,|\,B) \implies$ SI.
    \item Si $a = 3$: $\operatorname{Rg}(A) = \operatorname{Rg}(A\,|\,B) = 1 \neq 3 = $ nº incógnitas $\implies$ SCI.
\end{itemize}

En el caso $a=3$, el sistema equivalente es:

\[
\begin{cases}
x + 4y + 3z = 2
\end{cases}
\]

\[
\text{Incógnita principal: } x, \quad \text{Parámetros: } y, z.
\]
\[
x = 2 - 4y - 3z
\]

$y = \lambda_1$, $z = \lambda_2$. Entonces las soluciones son:

\[
\boxed{
x = 2 + 3\lambda_1 + 4\lambda_2, \quad y = \lambda_1, \quad z = \lambda_2, \qquad \lambda_1, \lambda_2 \in \mathbb{Z}_7
}
\]

\end{ejemplo}

\begin{procedimiento}[title=Sistema homogéneo]
Un sistema de ecuaciones lineales es homogéneo si todos los términos independientes valen $0$. Los sistemas homogéneos son siempre compatibles, ya que admiten al menos la solución trivial:
\[
x_1 = x_2 = \dots = x_n = 0.
\]
\end{procedimiento}

\section{Matrices Regulares}

\begin{procedimiento}[title=Matriz Regular]
Una matriz regular o invertible es una matriz cuadrada $A \in \mathfrak{M}_n(\mathbb{K})$ para la cual existe una única matriz $B \in \mathfrak{M}_n(\mathbb{K})$ tal que
\[
A \cdot B = I_n = B \cdot A,
\]
llamada la inversa de $A$: $B = A^{-1}$.
Además, $B$ es regular y $B^{-1} = A$.
\end{procedimiento}

\begin{procedimiento}[title=Propiedades de matrices regulares]
Sean $A, B \in \mathfrak{M}_n(\mathbb{K})$:
\begin{itemize}
    \item $I_n$ es regular y $(I_n)^{-1} = I_n$.
    \item $A$ regular $\implies$ $A^{-1}$ regular y $(A^{-1})^{-1} = A$.
    \item $A, B$ regulares $\implies$ $A \cdot B$ regular y $(A \cdot B)^{-1} = B^{-1} \cdot A^{-1}$.
    \item $A$ regular, se define $A^0 = I_n$.
    \item $(A^{-1})^n = (A^n)^{-1}$.
    \item $A$ regular $\iff \operatorname{Rg}(A) = n \iff H_f(A) = H_c(A) = I_n$.
\end{itemize}
\end{procedimiento}

\begin{procedimiento}[title=Cálculo de la inversa de una matriz]
Sea $A \in \mathfrak{M}_n(\mathbb{K})$. Para calcular $A^{-1}$:
\[
(A \,|\, I_n) \sim_f \dots \sim_f (I_n \,|\, B)
\quad \Rightarrow \quad B = A^{-1}.
\]
\end{procedimiento}

\begin{ejemplo}[title=Ejemplo de cálculo de $A^{-1}$ en $\mathbb{R}$]
Sea
\[
A =
\begin{pmatrix}
1 & 2 \\
3 & 5
\end{pmatrix} \in \mathfrak{M}_2(\mathbb{R}).
\]

\[
(A\,|\,I_2) =
\begin{pmatrix}
1 & 2 & \Big| & 1 & 0 \\
3 & 5 & \Big| & 0 & 1
\end{pmatrix}
\overset{\substack{F_2 - 3 F_1 \\{}}}{\sim_f}
\begin{pmatrix}
1 & 2 & \Big| & 1 & 0 \\
0 & -1 & \Big| & -3 & 1
\end{pmatrix}
\]

Vemos que $\operatorname{Rg}(A) = \text{número de filas no nulas en } 
\bigl(\!\begin{smallmatrix} 1 & 2 \\ 0 & -1 \end{smallmatrix}\!\bigr) 
= 2 = n \implies A \text{ es regular.}$

\[
\begin{pmatrix}
1 & 2 & \Big| & 1 & 0 \\
0 & -1 & \Big| & -3 & 1
\end{pmatrix}
\overset{\substack{F_1 + 2 F_2 \\ -F_2 \\{}}}{\sim_f}
\begin{pmatrix}
1 & 0 & \Big| & -5 & 2 \\
0 & 1 & \Big| & 3 & -1
\end{pmatrix}
\]

\[
\Rightarrow A^{-1} =
\begin{pmatrix}
-5 & 2 \\
3 & -1
\end{pmatrix}.
\]
\end{ejemplo}

\begin{ejemplo}[title=Ejemplo de matriz no regular en $\mathbb{Z}_7$]
Sea
\[
A =
\begin{pmatrix}
2 & 3 & 5 \\
4 & 2 & 1 \\
5 & 1 & 4
\end{pmatrix} \in \mathfrak{M}_3(\mathbb{Z}_7)
\]

\[
(A\,|\,I_3) =
\begin{pmatrix}
2 & 3 & 5 & \Big| & 1 & 0 & 0 \\
4 & 2 & 1 & \Big| & 0 & 1 & 0 \\
5 & 1 & 4 & \Big| & 0 & 0 & 1
\end{pmatrix}
\overset{\substack{F_2 - 2 F_1 \\ F_3 + F_1 \\{}}}{\sim_f}
\begin{pmatrix}
2 & 3 & 5 & \Big| & 1 & 0 & 0 \\
0 & 3 & 5 & \Big| & 5 & 1 & 0 \\
0 & 4 & 2 & \Big| & 1 & 0 & 1
\end{pmatrix}
\]
\[
\begin{pmatrix}
2 & 3 & 5 & \Big| & 1 & 0 & 0 \\
0 & 3 & 5 & \Big| & 5 & 1 & 0 \\
0 & 4 & 2 & \Big| & 1 & 0 & 1
\end{pmatrix}
\overset{\substack{F_3 + 2 F_2 \\{}}}{\sim_f}
\begin{pmatrix}
2 & 3 & 5 & \Big| & 1 & 0 & 0 \\
0 & 3 & 5 & \Big| & 5 & 1 & 0 \\
0 & 0 & 0 & \Big| & 6 & 1 & 1
\end{pmatrix}
\]

Vemos que $\operatorname{Rg}(A) = \text{número de filas no nulas en } 
\bigl(\!\begin{smallmatrix} 2 & 3 & 5 \\ 0 & 3 & 5 \\ 0 & 0 & 0 \end{smallmatrix}\!\bigr)
= 2 \neq n = 3 \implies A \text{ no es regular.}$
\end{ejemplo}

\section{Determinantes}

\begin{procedimiento}[title=Definición de determinante]
Sea 
\[
\mathfrak{M}(\mathbb{K}) = \bigcup_{\substack{n \in \mathbb{N} \\ n \ge 1}} \mathfrak{M}_n(\mathbb{K})
\]
El determinante de $A \in \mathfrak{M}_n(\mathbb{K})$ es una aplicación:
\[
\begin{aligned}
\mathfrak{M}(\mathbb{K}) &\longrightarrow \mathbb{K} \\
A &\longmapsto |A|
\end{aligned}
\]
\end{procedimiento}

\begin{procedimiento}[title=Cálculo de determinantes]
El determinante de $A \in \mathfrak{M}_n(\mathbb{K})$ es una suma y resta de $n!$ términos y se calcula mediante la \textbf{Regla de Sarrus} para los siguientes $n$:

\vspace{0.5em}

\textbf{$n = 1$}: 
\[
|a| = a
\]
\textbf{$n = 2$}: 
\[
\left|\begin{array}{cc}
a_{11} & a_{12} \\
a_{21} & a_{22}
\end{array}\right|= a_{11} a_{22} - a_{21} a_{12}
\]
\textbf{$n = 3$}:
\[
\left|\begin{array}{ccc}
a_{11} & a_{12} & a_{13} \\
a_{21} & a_{22} & a_{23} \\
a_{31} & a_{32} & a_{33}
\end{array}\right|
=
\begin{aligned}[t]
& a_{11}a_{22}a_{33} + a_{21}a_{32}a_{13} + a_{12}a_{23}a_{31} \\
& - a_{31}a_{22}a_{13} - a_{32}a_{23}a_{11} - a_{21}a_{12}a_{33}
\end{aligned}
\]


\end{procedimiento}

\begin{ejemplo}[title=Ejemplos de determinantes]
    \[
    A = (5) \implies |A| = 5
    \]

    \[
    A =
    \begin{pmatrix}
    3 & 7 \\
    2 & 5
    \end{pmatrix}, \quad
    |A| = 3\cdot 5 - 2\cdot 7 = 15 - 14 = 1
    \]

    \[
    A =
    \begin{pmatrix}
    1 & 2 & 3 \\
    0 & 4 & 5 \\
    1 & 0 & 6
    \end{pmatrix}, \quad
    |A| =
    \begin{aligned}[t]
    & 1\cdot 4 \cdot 6 + 0\cdot 0 \cdot 3 + 2\cdot 5 \cdot 1 \\
    & - 1\cdot 4 \cdot 3 - 0\cdot 5 \cdot 1 - 0\cdot 2 \cdot 6
    \end{aligned}
    = 24 + 10 - 12 = 22
    \]
\end{ejemplo}


\begin{procedimiento}[title=Propiedades del determinante]
\begin{enumerate}
    \item $|I_n| = 1$ para todo $n \ge 1$.
    \item Es lineal respecto a las filas y las columnas:  
    \[
    \left|
    \begin{matrix}
    F_1 \\
    \vdots \\
    \lambda F_i \\
    \vdots \\
    F_n
    \end{matrix}
    \right| = \lambda 
    \left|
    \begin{matrix}
    F_1 \\
    \vdots \\
    F_i \\
    \vdots \\
    F_n
    \end{matrix}
    \right|, 
    \quad
    \left|
    \begin{matrix}
    F_1 \\
    \vdots \\
    F_i + F_i' \\
    \vdots \\
    F_n
    \end{matrix}
    \right| =
    \left|
    \begin{matrix}
    F_1 \\
    \vdots \\
    F_i \\
    \vdots \\
    F_n
    \end{matrix}
    \right| +
    \left|
    \begin{matrix}
    F_1 \\
    \vdots \\
    F_i' \\
    \vdots \\
    F_n
    \end{matrix}
    \right|
    \]
    \[
    |C_1 \;\cdots\; \lambda C_i \;\cdots\; C_n| = \lambda \, |C_1 \;\cdots\; C_i \;\cdots\; C_n|, 
    \]
    \[
    |C_1 \;\cdots\; C_i + C_i' \;\cdots\; C_n| = |C_1 \;\cdots\; C_i \;\cdots\; C_n| + |C_1 \;\cdots\; C_i' \;\cdots\; C_n|
    \]
    \item $|A| = |A^T|$.
    \item Intercambiar dos líneas cambia el signo del determinante.
    \item |A| = 0 si y sólo si las líneas de A son linealmente dependientes:
    \begin{itemize}
        \item Si hay una línea nula, $|A| = 0$.
        \item Si en $A$ hay dos líneas iguales, $|A| = 0$.
        \item Si una línea es combinación lineal de otras, $|A| = 0$.
        \item Si a una línea se le suma una combinación lineal de otras, el determinante no varía.
    \end{itemize}

\end{enumerate}
\end{procedimiento}

\begin{procedimiento}[title=Combinación lineal de líneas]
Sea $A \in \mathfrak{M}_n(\mathbb{K})$ y $L_1, L_2, \dots, L_n$ las filas o columnas de $A$.  
Una línea $L$ es combinación lineal de las líneas $L_{i_1}, L_{i_2}, \dots, L_{i_r}$ si existen 
$\lambda_1, \lambda_2, \dots, \lambda_r \in \mathbb{K}$ tales que
\[
L = \lambda_1 L_{i_1} + \lambda_2 L_{i_2} + \dots + \lambda_r L_{i_r}.
\]
\end{procedimiento}

\begin{ejemplo}[title=Ejemplo de combinación lineal de líneas]
Sea
\[
A =
\begin{pmatrix}
1 & 2 & 3 \\
2 & 4 & 6 \\
3 & 5 & 7
\end{pmatrix} \in \mathfrak{M}_3(\mathbb{R}).
\]
La segunda fila es $2$ veces la primera fila:
\[
F_2 = 2 F_1,
\]
por lo que $F_2$ es combinación lineal de $F_1$.
\end{ejemplo}

\begin{procedimiento}[title=(In)dependencia lineal]
Las líneas $L_{i_1}, L_{i_2}, \dots, L_{i_r}$ son \textbf{linealmente dependientes} si existen 
$\lambda_1, \dots, \lambda_r \in \mathbb{K}$, no todos cero, tales que
\[
0 = \lambda_1 L_{i_1} + \dots + \lambda_r L_{i_r}.
\]

Las líneas $L_{i_1}, L_{i_2}, \dots, L_{i_r}$ son \textbf{linealmente independientes} si no son linealmente dependientes, es decir, si la única combinación lineal que da cero es
\[
0 = \lambda_1 L_{i_1} + \dots + \lambda_r L_{i_r} \implies \lambda_1 = \dots = \lambda_r = 0.
\]
\end{procedimiento}

\begin{ejemplo}[title=Ejemplo de líneas linealmente dependientes]
Sea
\[
A =
\begin{pmatrix}
1 & 2 & 3 \\
2 & 4 & 6 \\
3 & 5 & 7
\end{pmatrix} \in \mathfrak{M}_3(\mathbb{R}).
\]
Aquí, la segunda fila es $F_2 = 2 F_1$, y la tercera fila puede expresarse como combinación lineal de las dos primeras: $F_3 = -1 F_1 + 1 F_2$.  
Por lo tanto, las filas son linealmente dependientes.
\end{ejemplo}

\begin{ejemplo}[title=Ejemplo de líneas linealmente independientes]
Sea
\[
B =
\begin{pmatrix}
1 & 0 & 0 \\
0 & 1 & 0 \\
0 & 0 & 1
\end{pmatrix} \in \mathfrak{M}_3(\mathbb{R}).
\]
Las filas de $B$ no pueden expresarse como combinación lineal unas de otras salvo con todos los coeficientes cero.  
Por lo tanto, las filas son linealmente independientes.
\end{ejemplo}

\subsection{Matriz Adjunta}

\begin{procedimiento}[title=Matriz Adjunta]
La matriz adjunta de $A = (a_{ij}) \in \mathfrak{M}_n(\mathbb{K})$ es la matriz $\operatorname{Adj}(A) = \overline{A} = (\alpha_{ij}) \in \mathfrak{M}_n(\mathbb{K})$ donde
\[
\alpha_{ij} = (-1)^{i+j} \, |A_{ij}|,
\]
y $A_{ij}$ es la submatriz de $A$ que resulta de suprimir la fila $i$ y la columna $j$.
\end{procedimiento}

\begin{ejemplo}[title=Ejemplo de cálculo de la matriz adjunta]
Sea
\[
A =
\begin{pmatrix}
1 & 1 & 2 \\
2 & 4 & 5 \\
4 & 3 & 6
\end{pmatrix} \in \mathfrak{M}_3(\mathbb{Z}_7).
\]

Calculamos los cofactores $\alpha_{ij}$:

\[
\begin{aligned}
\alpha_{11} &= (-1)^{1+1} \, \left|
\begin{matrix} 4 & 5 \\ 3 & 6 \end{matrix}\right| = 4\cdot 6 - 5\cdot 3 = 24-15 = 9 \equiv 2, \\
\alpha_{12} &= (-1)^{1+2} \, \left|
\begin{matrix} 2 & 5 \\ 4 & 6 \end{matrix}\right| = -(2\cdot 6 - 5\cdot 4) = -(12-20) = 8 \equiv 1, \\
\alpha_{13} &= (-1)^{1+3} \, \left|
\begin{matrix} 2 & 4 \\ 4 & 3 \end{matrix}\right| = 2\cdot 3 - 4\cdot 4 = 6-16 = -10 \equiv 4, \\
\alpha_{21} &= (-1)^{2+1} \, \left|
\begin{matrix} 1 & 2 \\ 3 & 6 \end{matrix}\right| = -(1\cdot 6 - 2\cdot 3) = -(6-6)=0, \\
\alpha_{22} &= (-1)^{2+2} \, \left|
\begin{matrix} 1 & 2 \\ 4 & 6 \end{matrix}\right| = 1\cdot 6 - 2\cdot 4 = 6-8=-2 \equiv 5, \\
\alpha_{23} &= (-1)^{2+3} \, \left|
\begin{matrix} 1 & 1 \\ 4 & 3 \end{matrix}\right| = -(1\cdot 3 -1\cdot 4)=-(3-4)=1, \\
\alpha_{31} &= (-1)^{3+1} \, \left|
\begin{matrix} 1 & 2 \\ 4 & 5 \end{matrix}\right| = 1\cdot 5 - 2\cdot 4 = 5-8=-3 \equiv 4, \\
\alpha_{32} &= (-1)^{3+2} \, \left|
\begin{matrix} 1 & 2 \\ 2 & 5 \end{matrix}\right| = -(1\cdot 5 - 2\cdot 2)=-(5-4)=-1 \equiv 6, \\
\alpha_{33} &= (-1)^{3+3} \, \left|
\begin{matrix} 1 & 1 \\ 2 & 4 \end{matrix}\right| = 1\cdot 4 -1\cdot 2 = 2.
\end{aligned}
\]

Por lo tanto,
\[
\overline{A} =
\begin{pmatrix}
2 & 1 & 4 \\
0 & 5 & 1 \\
4 & 6 & 2
\end{pmatrix} \in \mathfrak{M}_3(\mathbb{Z}_7).
\]
\end{ejemplo}

\begin{procedimiento}[title=Propiedad de la matriz adjunta]
Para toda $A \in \mathfrak{M}_n(\mathbb{K})$ se cumple
\[
A \cdot (\overline{A})^T = |A| \, I_n = (\overline{A})^T \cdot A.
\]
\end{procedimiento}

\subsection{Propiedades consecuencia}

\subsubsection{Desarrollos de Laplace}

\begin{procedimiento}[title=Desarrollos de Laplace]
Para $A \in \mathfrak{M}_n(\mathbb{K})$:
\begin{itemize}
    \item Desarrollo por la fila $i$:
    \[
    |A| = \sum_{j=1}^n a_{ij} \, \alpha_{ij}, \quad \forall i \in \{1,\dots,n\}.
    \]
    \item Desarrollo por la columna $j$:
    \[
    |A| = \sum_{i=1}^n a_{ij} \, \alpha_{ij}, \quad \forall j \in \{1,\dots,n\}.
    \]
\end{itemize}
\end{procedimiento}

\begin{ejemplo}[title=Ejemplo de desarrollos de Laplace]
Sea
\[
A =
\begin{pmatrix}
1 & 1 & 2 \\
2 & 4 & 5 \\
4 & 3 & 6
\end{pmatrix} \in \mathfrak{M}_3(\mathbb{Z}_7). \qquad
\overline{A} =
\begin{pmatrix}
2 & 1 & 4 \\
0 & 5 & 1 \\
4 & 6 & 2
\end{pmatrix}.
\]

El determinante es
\[
|A| = 4 \in \mathbb{Z}_7.
\]

Desarrollo de Laplace por la segunda fila:
\[
|A| = 2\cdot 0 + 4\cdot 5 + 5\cdot 1 = 0 + 20 + 5 \equiv 4.
\]

Desarrollo de Laplace por la tercera columna:
\[
|A| = 2\cdot 4 + 5\cdot 1 + 6\cdot 2 = 8 + 5 + 12 \equiv 4.
\]

Se cumple la propiedad, ambos desarrollos coinciden.
\end{ejemplo}

\subsubsection{Caracterización de matrices regulares y fórmula para la inversa}

\begin{procedimiento}[title=Propiedad y fórmula para la inversa]
\begin{itemize}
    \item Para todas $A,B \in \mathfrak{M}_n(\mathbb{K})$,
    \[
    |A B| = |A| \cdot |B|.
    \]

    \item Una matriz $A \in \mathfrak{M}_n(\mathbb{K})$ es \textbf{regular} si y sólo si \(|A| \neq 0.\)
    
    En tal caso,
    \[
    A^{-1} = \frac{1}{|A|} \, (\overline{A})^{T}.
    \]
\end{itemize}
\end{procedimiento}

\begin{ejemplo}[title=Ejemplo de cálculo de la inversa con la adjunta]
Sea
\[
A =
\begin{pmatrix}
1 & 1 & 2 \\
2 & 4 & 5 \\
4 & 3 & 6
\end{pmatrix}
\in \mathfrak{M}_3(\mathbb{Z}_7),
\qquad
\overline{A} =
\begin{pmatrix}
2 & 1 & 4 \\
0 & 5 & 1 \\
4 & 6 & 2
\end{pmatrix}.
\]

Sabemos que
\[
|A| = 4 \in \mathbb{Z}_7.
\]

Entonces, usando la fórmula
\[
A^{-1} = \frac{1}{|A|}\, (\overline{A})^T,
\]
primero calculamos el inverso multiplicativo de $4$ en $\mathbb{Z}_7$:
\[
4^{-1} = 2.
\]

Por tanto,
\[
A^{-1}
= 2 \cdot
\begin{pmatrix}
2 & 0 & 4 \\
1 & 5 & 6 \\
4 & 1 & 2
\end{pmatrix}
=
\begin{pmatrix}
4 & 0 & 1 \\
2 & 3 & 5 \\
1 & 2 & 4
\end{pmatrix}.
\]

\end{ejemplo}

\subsubsection{Otros significados de $\operatorname{Rg}(A)$}

\begin{procedimiento}[title=Rango como mayor orden de submatrices regulares]
Sea $A \in \mathfrak{M}_{m \times n}(\mathbb{K})$.  
$\operatorname{Rg}(A)$ es el mayor orden de las submatrices regulares de $A$ (submatrices cuyo determinante es no nulo).

\medskip
\textbf{Observación:}  
Si $A \in \mathfrak{M}_n(\mathbb{K})$ es regular, entonces $\operatorname{Rg}(A) = n$.

\end{procedimiento}

\begin{ejemplo}
Sea
\[
A =
\begin{pmatrix}
1 & 1 & 2 \\
2 & 4 & 5 \\
4 & 3 & 6
\end{pmatrix}.
\]
Entonces
\[
|A| =
\begin{vmatrix}
1 & 1 & 2 \\
2 & 4 & 5 \\
4 & 3 & 6
\end{vmatrix}
= 4 \neq 0.
\]
Por tanto, $A$ es regular y
\[
\operatorname{Rg}(A)=3.
\]
\end{ejemplo}

\begin{procedimiento}[title=Rango como mayor número de líneas linealmente independientes]
Sea $A \in \mathfrak{M}_{m \times n}(\mathbb{K})$. $\operatorname{Rg}(A)$ es el \textbf{mayor número de lineas linealmente independientes} de $A$.

\end{procedimiento}

\begin{ejemplo}
Sea $A \in \mathfrak{M}_{3 \times 4}(\mathbb{Z}_7)$:
\[
A=
\begin{pmatrix}
3 & 5 & 1 & 2 \\
2 & 4 & 1 & 6 \\
4 & 3 & 3 & 5
\end{pmatrix}
\overset{\substack{F_3+F_1 \\ 5F_1\\{}}}{\sim_f}
\begin{pmatrix}
1 & 4 & 5 & 3 \\
2 & 4 & 1 & 6 \\
0 & 1 & 4 & 0
\end{pmatrix}
\overset{F_2-2F_1}{\sim_f}
\begin{pmatrix}
1 & 4 & 5 & 3 \\
0 & 3 & 5 & 0 \\
0 & 1 & 4 & 0
\end{pmatrix}
\]

\[
\begin{pmatrix}
1 & 4 & 5 & 3 \\
0 & 3 & 5 & 0 \\
0 & 1 & 4 & 0
\end{pmatrix}
\overset{F_2-2F_3}{\sim_f}
\begin{pmatrix}
1 & 4 & 5 & 3 \\
0 & 0 & 0 & 0 \\
0 & 1 & 4 & 0
\end{pmatrix}
\overset{F_2 \leftrightarrow F_3}{\sim_f}
\begin{pmatrix}
1 & 4 & 5 & 3 \\
0 & 1 & 4 & 0 \\
0 & 0 & 0 & 0
\end{pmatrix}.
\]

\vspace{0.5em}

Como las dos primeras filas son no nulas:  
\[
\operatorname{Rg}(A) = 2.
\]

\medskip
\textbf{Verificación por determinantes:}

\[
\begin{vmatrix}
3 & 5 \\
2 & 4
\end{vmatrix}
= 3\cdot 4 - 2\cdot 5 = 12 - 10 = 2 \neq 0 \;\Rightarrow\; \operatorname{Rg}(A)\ge 2.
\]

\[
\begin{vmatrix}
3 & 5 & 1 \\
2 & 4 & 1 \\
4 & 3 & 3
\end{vmatrix}
= 0
\quad
(\text{la tercera columna es combinación lineal de la primera y segunda}).
\]

\[
\begin{vmatrix}
3 & 5 & 2 \\
2 & 4 & 6 \\
4 & 3 & 5
\end{vmatrix}
= 0
\quad
(\text{la cuarta columna también es combinación de la primera y segunda}).
\]

Por tanto,
\[
\operatorname{Rg}(A)=2.
\]

\textbf{Verificación en las filas:}

Las filas 1 y 2 son linealmente independientes.

La fila 3 es combinación lineal de las dos primeras.

\vspace{0.5em}

Se sigue cumpliendo:

\[
\operatorname{Rg}(A)=2.
\]
\end{ejemplo}

\subsubsection{La fórmula de Cramer}

\begin{procedimiento}[title=Sistemas de Cramer]
Un sistema de ecuaciones lineales es un \textbf{sistema de Cramer} si se cumplen:
\begin{itemize}
    \item el número de ecuaciones coincide con el número de incógnitas.
    \item la matriz de coeficientes $A$ es regular.
\end{itemize}

En tal caso, el sistema es \textbf{SCD} y su única solución es
\[
X = A^{-1} B.
\]

\end{procedimiento}

\begin{ejemplo}
En $\mathbb{Z}_7$:
\[
\left\{
\begin{aligned}
x + y + 2z &= 3,\\
2x + 4y + 5z &= 4,\\
4x + 3y + 6z &= 1.
\end{aligned}
\right.
\]

La matriz ampliada es
\[
(A\,|\,B)=
\begin{pmatrix}
1 & 1 & 2 \,\,\Big|\,\, 3\\
2 & 4 & 5 \,\,\Big|\,\, 4\\
4 & 3 & 6 \,\,\Big|\,\, 1
\end{pmatrix}.
\]

Como $|A| = 4 \neq 0$, la matriz es regular $\implies$ es un sistema de Cramer.

La única solución es
\[
X=
\begin{pmatrix}
x\\y\\z
\end{pmatrix}
=
A^{-1}\cdot B
=
\begin{pmatrix}
4 & 0 & 1\\
2 & 3 & 5\\
1 & 2 & 4
\end{pmatrix}
\begin{pmatrix}
3\\4\\1
\end{pmatrix}
=
\begin{pmatrix}
6\\2\\1
\end{pmatrix}.
\]

Por tanto,  
\[
x = 6,\qquad y = 2,\qquad z = 1.
\]

\end{ejemplo}

\begin{procedimiento}[title=Fórmula de Cramer]
Sea $A X = B$ un sistema de Cramer con $A \in \mathfrak{M}_n(\mathbb{K})$.
La solución única viene dada por
\[
x_i = \frac{|M_i|}{|A|}, \qquad i = 1,2,\dots,n,
\]
donde $M_i$ se obtiene sustituyendo en $A$ la $i$-ésima columna por la matriz $B$.
\end{procedimiento}

\begin{ejemplo}
En $\mathbb{Z}_7$:
\[
\left\{
\begin{aligned}
x + y + 2z &= 3,\\
2x + 4y + 5z &= 4,\\
4x + 3y + 6z &= 1.
\end{aligned}
\right.
\]

\vspace{0.5em}

\[
(A\,|\,B)=
\begin{pmatrix}
1 & 1 & 2 \,\,\Big|\,\, 3\\
2 & 4 & 5 \,\,\Big|\,\, 4\\
4 & 3 & 6 \,\,\Big|\,\, 1
\end{pmatrix},
\qquad
|A| = 4 \neq 0.
\]

$\implies$ A es regular. Por la fórmula de Cramer:

\[
x = \frac{1}{4}\cdot
\begin{vmatrix}
3 & 1 & 2\\
4 & 4 & 5\\
1 & 3 & 6
\end{vmatrix}
= 2 \cdot 3 = 6.
\]

\[
y = \frac{1}{4}\cdot
\begin{vmatrix}
1 & 3 & 2\\
2 & 4 & 5\\
4 & 1 & 6
\end{vmatrix}
= 2 \cdot 1 = 2.
\]

\[
z = \frac{1}{4}\cdot
\begin{vmatrix}
1 & 1 & 3\\
2 & 4 & 4\\
4 & 3 & 1
\end{vmatrix}
= 2 \cdot 4 = 1.
\]

\end{ejemplo}

\section{Conceptos Adicionales}

\subsection{Traza}

\begin{procedimiento}[title=Traza]
La traza $\operatorname{tr}(A)$ de una matriz cuadrada $A$ es la suma de los elementos de su diagonal principal.
\end{procedimiento}

\begin{ejemplo}[title=Ejemplo de traza]
Sea
\[
A=
\begin{pmatrix}
\boxed{3} & 2 & 1 \\
4 & \boxed{5} & 0 \\
7 & 1 & \boxed{6}
\end{pmatrix}.
\]
Entonces
\[
\operatorname{tr}(A)=3+5+6=14.
\]
\end{ejemplo}

\begin{procedimiento}[title=Propiedades de la traza]
\begin{itemize}
    \item $\operatorname{tr}(A+B)=\operatorname{tr}(A)+\operatorname{tr}(B)$
    \item $\operatorname{tr}(\lambda A)=\lambda\,\operatorname{tr}(A)$
    \item $\operatorname{tr}(AB)=\operatorname{tr}(BA)$
\end{itemize}
\end{procedimiento}

% --------------------------------------------------------------------

\subsection{Matrices estocásticas o de probabilidad}

\begin{procedimiento}[title=Matriz estocástica o de probabilidad]
Una matriz cuadrada $A=(a_{ij})\in M_n(\mathbb{R})$ es estocástica o de probabilidad si cumple:
\begin{itemize}
    \item $a_{ij}\ge 0$ para todo $i,j$.
    \item Para cada columna, la suma de sus elementos es $1$.
\end{itemize}
\end{procedimiento}

\begin{ejemplo}[title=Ejemplo de matriz estocástica]
Sea
\[
A=
\begin{pmatrix}
0\text{.}5 & 0   & 0\text{.}25 \\
0\text{.}5 & 0\text{.}5 & 0\text{.}75 \\
0   & 0\text{.}5 & 0
\end{pmatrix}.
\]
Las columnas suman:
\[
0\text{.}5+0\text{.}5+0=1,\qquad
0+0\text{.}5+0\text{.}5=1,\qquad
0\text{.}25+0\text{.}75+0=1.
\]
Por tanto, $A$ es estocástica.
\end{ejemplo}

\begin{procedimiento}[title=Matriz doblemente estocástica]
Una matriz estocástica es doblemente estocástica si además cumple:
\begin{itemize}
    \item Para cada fila, la suma de sus elementos es $1$.
\end{itemize}
\end{procedimiento}

\begin{ejemplo}[title=Ejemplo de matriz doblemente estocástica]
Sea
\[
B=
\begin{pmatrix}
0\text{.}1 & 0\text{.}5 & 0\text{.}4 \\
0\text{.}6 & 0\text{.}2 & 0\text{.}2 \\
0\text{.}3 & 0\text{.}3 & 0\text{.}4
\end{pmatrix}.
\]
Cada columna suma $1$ y cada fila también suma $1$, por lo que $B$ es doblemente estocástica.
\end{ejemplo}

\end{document}
